\section{Numerical illustrations}
\label{se:numerical.experiments}
We provide in this section illustrations of the frequency and time focus functions introduced in the core of the paper. We stress that these do not intend to address specific applied problems, but simply to show  that such focus functions can indeed be designed and achieve well targeted goals.

Stepping from continuous time functions to discrete signals requires choosing a discretization scheme. Our approach here was to limit ourselves to uniform, frequency or scale independent, time sampling. In other words, we stick to very redundant time-frequency/scale transforms, and do not address discretization issues such as the ones developed in classical frame theory, which we consider beyond the scope of this paper. 

\subsection{Illustration of time focus}
\label{sec:app.time.foc.example}
We illustrate the time focus effect using a simple example of time focus function, applied to a real audio signal. For the sake of simplicity, we take $\gamma(\omega)=\omega$ for all $\omega\in\R$. Given some signal $f\in L^2(\R)$, we denote by $V f = M^\tau_{\sigma_\text{ref}}f$ the transform of $f$, with a focus function uniformly equal to a reference scale $\sigma_\text{ref}$, and define
\begin{equation}
\label{eq_time_focus_1}
\sigma_f^\tau(t) = A \int \omega^{n} |V f(t,\omega)|\,d\omega + B\ ,
\end{equation}
where $n\in\N$ is a fixed integer, and $A,B>0$ are real constants that can be adjusted so that for all $t$,
\[
1\le \sigma_f^\tau (t)\le\sigma_{\text{max}}\ ,
\]
for some prescribed maximal focus $\sigma_{\text{max}}$.

\begin{figure}[h!]
    \centering
    \includegraphics[scale=.5]{casta_signal}
    \caption{Castanet signal and corresponding time focus function, defined by Equation~\eqref{eq_time_focus_1}}.
    \label{fi:casta.signal}
\end{figure}
%We notice that, if the window $h$ belongs to the Feichtinger algebra $S_0(\R)$, the existence of the focus function is guaranteed as soon as $f$ belongs to the modulation space $f\in\mathcal{M}_{w}^{2,\infty}(\R)$, with weight function $w(t,\omega) = \sqrt{1+\omega^{2n}}$.

\medskip

We display in Fig.~\ref{fi:casta.signal} a 3.5 seconds excerpt from a castanet sound recording (from the SQAM assessment database~\cite{EBU1988SoundQA}), and the corresponding focus function estimated using Equation~\eqref{eq_time_focus_1}. The window $h$ was a $10$ milliseconds long truncated Gaussian window (to enforce compact support), and parameters were set to $n=1$, $\sigma_\text{ref}=1$ and $\sigma_\text{max}=5$. As can be seen, the transients are well detected. Fig.~\ref{fi:casta.spgrams} represents the spectrograms obtained with the unfocused transform ($\sigma(t)=1$ for all $t$), and the focused transform. The latter features sharper attacks, the invervals in between attacks being unchanged.
%We now revisit Example~\ref{ex:time_focus}, and illustrate the behavior of a time focus function similar to~\eqref{eq_def2_sigma_tau} on another toy example, consisting in colored noise with additional randomly located spikes with random amplitude. A realization of such a random signal (involving 10 random peaks, out of which some are not visible because they are hidden by the background noise) is plotted in  Figure~\ref{fi:time.foc.sig}. Its spectrogram (not shown here) exhibits a clear frequency decay (as a result of the choice of noise) away from the spike locations (which do not decay in the frequency domain). The time focus function used here is designed to capture such a behavior, which we illustrate below.

\begin{figure}
    \centering
    \hspace{-17.5mm}
    \includegraphics[width=.58\textwidth]{casta_unfocused}
    \hspace{-12mm}
    \includegraphics[width=.58\textwidth]{casta_focused}\hspace{-22mm}
    \caption{Log-spectrograms for the castanet signal plotted in Fig.~\ref{fi:casta.signal}. Left: unfocused; Right: focused, using the focus function defined in~\eqref{eq_time_focus_1} with parameters given in the text.}
    \label{fi:casta.spgrams}
\end{figure}

\begin{figure}[h!]
    \centering
    \hspace{-11mm}
    \includegraphics[scale=.45]{glock_signal}
    \hspace{-9mm}
    \includegraphics[scale=.45]{glock_signal_entropy}
    \hspace{-15mm}
    \caption{Glockenspiel signal (top) and corresponding time focus functions. Left: focus function as defined in Equation~\eqref{eq_time_focus_1}. Right: focus function as defined in Equation~\eqref{eq_time_focus_2}.}
    \label{fi:glock.signal}
\end{figure}

This example is quite an easy one, as the signal only contains transients. The same focus function performs worse on a slightly more complex signal, that features significant harmonic components together with transients. We display in Fig.~\ref{fi:glock.signal} a 3.5 seconds excerpt from a glockenspiel sound recording (available from the companion web site of~\cite{Jaillet2007time}), together with the corresponding focus function (bottom left-hand panel). As can be seen, the time focus function~\eqref{eq_time_focus_1} detects the attacks of notes, but the decay is much slower than it was for the castanet signal, and the focus effect on the resulting spectrogram (not shown here) is not satisfactory. In fact, the focus function in~\eqref{eq_time_focus_1} is indeed sensitive to transients, but also on the local energy of the signal. Increasing the value of $n$ does not seem to improve.


As an alternative, we display in the bottom right-hand panel of Fig.~\ref{fi:glock.signal} the focus function based upon the entropy of fixed-time slices of the reference spectrogram (suitably normalized to unit norm).
\begin{equation}
\label{eq_time_focus_2}
    \sigma_f(t) = - A \int |\widetilde V f(t,\omega)| \log |\widetilde V f(t,\omega)|\, d\omega\,+B\ ,
\end{equation}
where $\widetilde V f(t,\omega) = Vf(t,\omega)/\|V_f(t,\cdot)\|_{L^1(\R,d\omega)}$. Parameters $A$ and $B$ were again set to ensure $1\le\sigma(t)\le\sigma_\text{max}=5$.

The rationale is that slices that do not correspond to transient events exhibit a sparser behavior, and can therefore be expected to possess a small entropy. The right-hand panel of Fig.~\ref{fi:glock.signal} shows that the estimated focus function is indeed sensitive to transients, independently of the local amplitude (which is clear from the construction in~\eqref{eq_time_focus_2}).
The corresponding spectrograms are displayed in Fig.~\ref{fi:glock.spgrams.entropy}, from which a better focus effect can be seen on the transient attacks of the instrument. However, the sustained parts have lost their frequency resolution in parts of the signal featuring close transients (in the middle segment of the signal).

Entropy seems to be a valuable choice for building a time focus function. Let us nevertheless stress that the construction depends on several parameters, including the reference focus $\sigma_\text{ref}$ involved in the reference STFT $Vf$, and the maximal allowed value $\sigma_\text{max}$. One may also investigate extensions built upon Renyi entropies, which provide different measures of spreading in the time-frequency domain, as shown in~\cite{Jaillet2007time}.
\begin{figure}
    \centering
    \hspace{-17mm}
    \includegraphics[width=.58\textwidth]{glock_unfocused_entropy}
    \hspace{-12mm}
    \includegraphics[width=.58\textwidth]{glock_focused_entropy}
    \hspace{-22mm}
    \caption{Spectrograms for the glockenspiel signal and focus function plotted in the right hand panel of Fig.~\ref{fi:glock.signal}. Left: unfocused; Right: focused, using the entropy-based focus function defined in~\eqref{eq_time_focus_2}.}
    \label{fi:glock.spgrams.entropy}
\end{figure}

\begin{comment}
We display in Figure~\ref{fi:time.foc} the corresponding time focus function, as defined in~\eqref{eq_def2_sigma_tau}. The focus function presents clear maxima near 7 time instants, which turn out to coincide with the locations of the 7 neatly visible spikes (out of 10 generated spikes, the other 3 having smaller amplitude). The function was rescaled to achieve the prescribed minimal value of 1; again further post-processing is possible to achieve the support condition of the definition, we refrained from doing it as this would require an extra thresholding parameter, such a fine tuning being irrelevant at this stage.
\end{comment}


\subsection{Illustration of frequency focus}
\label{sec:app.freq.foc.example}
We now illustrate the behavior of the frequency-focused transform. Again, we will build a focus function using an entropy measure, based upon fixed-frequency slices of a standard continuous wavelet transform
\begin{equation}
\label{eq_freq_focus}
\sigma(u) = - A\int \left|\widetilde{W}f(t,u)\right| \log \left|\widetilde{W}f(t,u)\right| du\,+B\ ,    
\end{equation} 
where $\widetilde{W}f(t,u) = Wf(t,u)/\|Wf(\cdot,u)\|_{L^1(\R)}$ is a normalized continuous wavelet transform (equivalently a frequency-focused transform with focus function uniformly equal to $\sigma_\text{ref}=1$). We have chosen here the simplest choice $\gamma(u)=e^u$. Again, $A$ and $B$ are parameters which are adjusted so that $1\le\sigma(u)\le\sigma_\text{max}$, for all $u$ and for some prescribed maximal focus $\sigma_\text{max}$.
\begin{figure}[h!]
    \centering
    \hspace{-17mm}
    \includegraphics[scale=.45]{simul_sig_spgram}
    \hspace{-8mm}
    \includegraphics[scale=.45]{simul_freq_focus}
    \hspace{-20mm}
    \caption{Left: toy signal (top) and periodogram (bottom, loglog scale). Right: wavelet spectrum of the simulated signal and corresponding frequency focus function as defined in Equation~\eqref{eq_freq_focus}.}
    \label{fi:simul_freq_focus}
\end{figure}
We display in Fig.~\ref{fi:simul_freq_focus} the simulated signal and its periodogram (square modulus of Fourier transform) on the left, and its wavelet spectrum and the frequency focus. The  simulated signal is composed of the sum of four sine waves at different frequencies with equal amplitudes, randomly located spikes with random amplitudes (50 spikes) and Gaussian white noise. The wavelet spectrum is defined as the time-average of the continuous wavelet transform modulus displayed in Fig.~\ref{fi:simul.scgrams.entropy}, left panel.

Obviously, the frequency focus is insensitive to the different amplitudes of the four sine waves in the wavelet domain. The resulting effect is visible on the scalograms (modulus of time-scale transforms) on the right panel of Fig.~\ref{fi:simul.scgrams.entropy}, where the frequency resolution has clearly been increased for displaying the four sine waves, and is weakly changed elsewhere, in particular at smallest scales. It is also worth observing that the localization of spikes at small scales from wavelet maxima appears simpler, since these lines of maxima are less affected by the presence of the sine wave.
We didn't consider a real example for illustrating the frequency focus, since constant amplitude sine waves rarely appear in real signals. Most often, sine waves start at a given time and their amplitude decays with time, which is not accounted for by the simple criterion illustrated here. The latter could be adapted to be used inside time segments, after a prior time segmentation. Such an extension would hover require additional modeling work, and is beyond the scope of this paper.



\begin{figure}
    \centering
    \hspace{-17mm}
    \includegraphics[width=.58\textwidth]{simul_cwt}
    \hspace{-12mm}
    \includegraphics[width=.58\textwidth]{simul_fftft}
    \hspace{-22mm}
    \caption{Scalograms for the glockenspiel signal and frequency focus function plotted in Fig.~\ref{fi:simul_freq_focus}. Left: unfocused; Right: focused, using the entropy-based focus function defined in~\eqref{eq_freq_focus}.}
    \label{fi:simul.scgrams.entropy}
\end{figure}

\begin{comment}
We revisit here Example~\ref{ex:frequency_focus}, and illustrate the behavior of a frequency focus function similar to~\eqref{fo:ex.freq.foc} on a toy example of harmonic signal, involving sine waves with close frequencies in the low frequency range, and sine waves with better separated frequencies in the high frequency range. A spectrogram (square modulus of STFT) with suitably chosen analysis window is displayed in Figure~\ref{fi:freq.foc.spgram}. The spectrogram was computed using LTFATpy~\cite{ltfatpy}, the Python port of the Large Time-Frequency Analysis Toolbox~\cite{ltfat}. Clearly, the window choice for the STFT allows separation of sine waves in the high frequency part of the spectrogram, while in the low-frequency range these sine waves are mixed by the window, resulting in beating phenomenon (oscillations of the modulus due to interference between sine waves). Clearly enough, these two different behaviors result from a suitable choice of the STFT window scale.

\begin{figure}[h!]
    \centering
    \includegraphics[scale=.6]{freq_foc_spgram}
    \caption{Spectrogram of a harmonic signal involving sine waves with close frequencies in the low frequency range, and sine waves with better separated frequencies in the high frequency range.}
    \label{fi:freq.foc.spgram}
\end{figure}

\begin{figure}[h!]
    \centering
    \includegraphics[scale=.6]{freq_foc}
    \caption{Frequency focus function based upon a ratio $L^2/L^1$ norms, for the signal whose spectrogram is displayed in Figure~\ref{fi:freq.foc.spgram}}
    \label{fi:freq.foc}
\end{figure}
The prototypical frequency focus function introduced in~\eqref{fo:ex.freq.foc} turns out to be insufficient to produce the expected behavior, because of spurious oscillations (not shown here). The latter can be suppressed by appropriate low-pass filtering. After such filtering and normalization, on obtains a focus function as given in Figure~\ref{fi:freq.foc}. The latter can be suitably rescaled to match a prescribed range, and thresholded to achieve the support property assumed in its definition. Since it is not the objective of this paper to provide realistic choices for applications, we d'on't  address such a fine tuning here.

As can be seen on Figure~\ref{fi:freq.foc}, the frequency focus function takes lower values in the low-frequency part of the signal spectrum. This implies that the window in the frequency-focused transform will have larger scale, therefore smaller frequency bandwidth, which will provide the desired adaptation to the signal if the aim is to better separate sine waves.

\vspace{-5mm}
\begin{remark}
Other more explicit illustrations can also be obtained since all calculations rely on Gaussian integrals which are tractable. In our opinion this does not bring more insight than the numerical examples in the Appendix, we then refrain from following this avenue.
\end{remark}
\end{comment}